% Part: first-order-logic
% Chapter: natural-deduction
% Section: identity

\documentclass[../../../include/open-logic-section]{subfiles}

\begin{document}

\olfileid{fol}{ntd}{ide}

\olsection{\usetoken{S}{identity}తో \usetoken{P}{derivation}}

!!{identity}తో కూడిన !!^{derivation}sకు అదనపు నిగమన నియమాలు అవసరం.

\begin{defish}
\AxiomC{}
\RightLabel{\Intro{\eq}}
\UnaryInfC{$\eq[t][t]$}
\DisplayProof
\hfill
\begin{tabular}{r}
\AxiomC{$\eq[t_1][t_2]$}
\AxiomC{$!A(t_1)$}
\RightLabel{\Elim{\eq}}
\BinaryInfC{$!A(t_2)$}
\DisplayProof
\\[3ex]
\AxiomC{$\eq[t_1][t_2]$}
\AxiomC{$!A(t_2)$}
\RightLabel{\Elim{\eq}}
\BinaryInfC{$!A(t_1)$}
\DisplayProof
\end{tabular}
\end{defish}

పై నియమాల్లో $t$, $t_1$, $t_2$లు సంవృత పదాలు. \Intro{\eq} నియమం
ఎలాంటి పరికల్పనలూ లేకుండానే $\eq[t][t]$ రూపంలోని ఏ తాదాత్మ్య
ప్రకటననైనా నేరుగా !!{derive} చేయనిస్తుంది.

\begin{ex}
$s$, $t$లు సంవృత పదాలైతే $!A(s), \eq[s][t] \Proves !A(t)$:
\begin{prooftree}
\AxiomC{$\eq[s][t]$}
\AxiomC{$!A(s)$}
\RightLabel{$\Elim{\eq}$}
\BinaryInfC{$!A(t)$}
\end{prooftree}
ఇది “తాదాత్మ్యాల ప్రతిస్థాపనీయతా సూత్రం” లేదా లైబ్నిజ్ నియమంగా
పరిచయం ఉండవచ్చు.
\end{ex}

\begin{prob}
$=$ సౌష్ఠవమైనదీ సంక్రామకమైనదీ అని నిరూపించండి; అంటే
$\lforall[x][\lforall[y][(\eq[x][y] \lif \eq[y][x])]]$,
$\lforall[x][\lforall[y][\lforall[z]((\eq[x][y] \land \eq[y][z])
\lif \eq[x][z])]]$ యొక్క !!{derivation}s ఇవ్వండి.
\end{prob}

\begin{ex}
!!{sentence}
\begin{align*}
& \lforall[x][\lforall[y][((!A(x) \land !A(y)) \lif \eq[x][y])]]
\intertext{ను ఈ !!{sentence} నుంచి !!{derive} చేస్తాం:}
& \lexists[x][\lforall[y][(!A(y) \lif \eq[y][x])]]
\end{align*}
!!{derivation}ను వెనుక నుంచి నిర్మిద్దాం:
\begin{prooftree}
\AxiomC{$\lexists[x][\lforall[y][(!A(y) \lif \eq[y][x])]]
\quad \Discharge{!A(a) \land !A(b)}{1}$}
\DeduceC{$\eq[a][b]$}
\DischargeRule{\Intro{\lif}}{1}
\UnaryInfC{$((!A(a) \land !A(b)) \lif \eq[a][b])$}
\RightLabel{\Intro{\lforall}}
\UnaryInfC{$\lforall[y][((!A(a) \land !A(y)) \lif \eq[a][y])]$}
\RightLabel{\Intro{\lforall}}
\UnaryInfC{$\lforall[x][\lforall[y][((!A(x) \land !A(y)) \lif \eq[x][y])]]$}
\end{prooftree}
ఇప్పుడు ప్రధాన పరికల్పనను వాడాలి: అది ఒక అస్తిత్వ !!{formula} కనుక
మధ్యంతర నిష్కర్ష $\eq[a][b]$ను !!{derive} చేయడానికి
\Elim{\lexists}ను వాడతాం.
\begin{prooftree}
\AxiomC{$\lexists[x][\lforall[y][(!A(y) \lif \eq[y][x])]]$}
\AxiomC{$\Discharge{\lforall[y][(!A(y) \lif \eq[y][c]])}{2}$}
\noLine
\UnaryInfC{$\Discharge{!A(a) \land !A(b)}{1}$}
\DeduceC{$\eq[a][b]$}
\DischargeRule{\Elim{\lexists}}{2}
\BinaryInfC{$\eq[a][b]$}
\DischargeRule{\Intro{\lif}}{1}
\UnaryInfC{$((!A(a) \land !A(b)) \lif \eq[a][b])$}
\RightLabel{\Intro{\lforall}}
\UnaryInfC{$\lforall[y][((!A(a) \land !A(y)) \lif \eq[a][y])]$}
\RightLabel{\Intro{\lforall}}
\UnaryInfC{$\lforall[x][\lforall[y][((!A(x) \land !A(y)) \lif \eq[x][y])]]$}
\end{prooftree}
కుడివైపు పైనున్న ఉప-!!{derivation}ను దాని పరికల్పనలతో
$\eq[a][c]$, $\eq[b][c]$ అని చూపి పూర్తి చేస్తాం. ఇందుకు రెండు వేరు
!!{derivation}s అవసరం. $\eq[a][c]$కు !!{derivation} కిందివిధంగా ఉంటుంది:
\begin{prooftree}
\AxiomC{$\Discharge{\lforall[y][(!A(y) \lif \eq[y][c]])}{2}$}
\RightLabel{\Elim{\lforall}}
\UnaryInfC{$!A(a) \lif \eq[a][c]$}
\AxiomC{$\Discharge{!A(a) \land !A(b)}{1}$}
\RightLabel{\Elim{\land}}
\UnaryInfC{$!A(a)$}
\RightLabel{\Elim{\lif}}
\BinaryInfC{$\eq[a][c]$}
\end{prooftree}
$\eq[a][c]$, $\eq[b][c]$ నుంచి \Elim{\eq} ద్వారా $\eq[a][b]$ను
!!{derive} చేస్తాం.
\end{ex}

\begin{prob}
కింది !!{formula}s యొక్క !!{derivation}s ఇవ్వండి:
\begin{enumerate}
\item $\lforall[x][\lforall[y][((\eq[x][y] \land !A(x)) \lif !A(y))]]$
\item $\lexists[x][!A(x)] \land \lforall[y][\lforall[z][((!A(y) \land
    !A(z)) \lif \eq[y][z])]] \lif \lexists[x][(!A(x) \land
  \lforall[y][(!A(y) \lif \eq[y][x])])]$
\end{enumerate}
\end{prob}

\end{document}
